\clearpage
\subsection{Cognitive Load — NASA Raw Task Load Index}
\label{sec:eval-rtlx}

Where SUS measured how usable the interfaces felt, the Raw Task
Load Index decomposes the participant's effort along six
interpretable axes (Mental Demand, Physical Demand, Temporal
Demand, Performance, Effort, and Frustration), each rated on the
standard NASA 21-point scale ($1$ = lowest demand, $21$ =
highest). Performance is reported in the positive-coded direction
(higher equals worse perceived performance) for internal
consistency with the other subscales. The overall R-TLX in this
study is the unweighted mean of the six subscales on the same
$1$--$21$ scale (no rescaling is applied). Full subscale
descriptions are in Appendix~\ref{app:sus-rtlx-items};
per-subscale distributions and box plots are in
Appendix~\ref{app:items-rtlx}. A short supplementary paragraph at
the end of this section reports the objective per-mode completion
time alongside the perceptual subscales, as a behavioral
counterpart to Temporal Demand.

Figure~\ref{fig:rtlx-aggregate} reports the composite R-TLX score
per mode against the scale midpoint of $10.5$, which marks the
boundary between the lower and upper halves of the $1$--$21$
range. All three modes sit comfortably in the lower half: Form
scores lowest with a composite of $5.50$, while Ask and Chat are
closely tied at $6.33$ each. As with SUS, Form is the lightest of
the three at the aggregate level, with Ask and Chat following
close behind by roughly $0.8$ scale points.

\begin{figure}[htbp]
\centering
\includegraphics[width=0.85\linewidth]{rtlx_aggregate_barplot.png}
\caption{Mean R-TLX composite workload score per mode ($n = 69$),
    unweighted average of the six subscales on the
$1$--$21$ scale. The dashed horizontal line marks the scale
midpoint ($10.5$), shown for visual reference.}
\label{fig:rtlx-aggregate}
\end{figure}

\begin{table}[htbp]
\centering
\caption{NASA R-TLX subscale mean (SD) by mode ($n = 69$ per mode),
$1$--$21$ scale. The Aggregated column averages the three per-mode means.}
\label{tab:rtlx-per-subscale}
\begin{tabular}{@{}l c c c c@{}}
\toprule
Subscale & Aggregated & Form & Ask & Chat \\
\midrule
Mental Demand      & 6.31 (4.75) & 5.77 (4.26) & 6.80 (4.63) & 6.36 (5.31) \\
Physical Demand    & 3.19 (3.40) & 2.77 (2.71) & 3.54 (3.92) & 3.26 (3.46) \\
Temporal Demand    & 4.89 (4.28) & 4.16 (3.82) & 4.97 (4.15) & 5.54 (4.78) \\
Effort             & 6.13 (4.57) & 5.68 (4.56) & 6.39 (4.41) & 6.32 (4.77) \\
Frustration        & 6.86 (5.24) & 5.88 (4.84) & 7.58 (5.34) & 7.13 (5.44) \\
Performance (pos.) & 8.93 (5.82) & 8.71 (5.48) & 8.68 (5.64) & 9.39 (6.36) \\
\bottomrule
\end{tabular}
\end{table}

\begin{figure}[htbp]
\centering
\includegraphics[width=0.45\linewidth]{rtlx_mental_demand.png}\hfill
\includegraphics[width=0.45\linewidth]{rtlx_physical_demand.png}\\[1ex]
\includegraphics[width=0.45\linewidth]{rtlx_temporal_demand.png}\hfill
\includegraphics[width=0.45\linewidth]{rtlx_effort.png}\\[1ex]
\includegraphics[width=0.45\linewidth]{rtlx_frustration.png}\hfill
\includegraphics[width=0.45\linewidth]{rtlx_performance.png}
\caption{Box plots of the six R-TLX subscales by mode
($n = 69$ per mode), one plot per subscale.}
\label{fig:rtlx-subquestions}
\end{figure}

Subscale by subscale, Table~\ref{tab:rtlx-per-subscale} reveals
three patterns in how participants rated workload.

Participants found Form somewhat less mentally demanding than Ask
or Chat (Form $5.77$, Ask $6.80$, Chat $6.36$), a gap of about
$0.6$ to $1.0$ points. The same small lead reappears on Effort
(Form $5.68$ vs Ask $6.39$, Chat $6.32$) and on Physical Demand
(Form $2.77$, Ask $3.54$, Chat $3.26$); the Physical Demand
ratings sit near the floor of the scale in all three modes, since
the ROPA task is not physically taxing in any of them.
The two subscales where Form's lead is sharpest are also the two
that reach statistical significance. Participants felt under
noticeably more time pressure in Chat than in Form (Form $4.16$,
Ask $4.97$, Chat $5.54$); the gap widens monotonically as LLM
support increases, the largest such increase across the six
subscales. Frustration shows the widest gap of any subscale: Ask
was rated the most frustrating ($7.58$), Chat close behind
($7.13$), and Form clearly the least ($5.88$). These are the
workload dimensions where LLM-driven interaction measurably costs
the user.
Performance is the one subscale where the Form lead disappears.
Participants rated their own success at the task as essentially
identical between Form ($8.71$) and Ask ($8.68$), with Chat
marginally worse ($9.39$, higher equals worse perceived
performance). Mode does not move how successfully participants
think they completed the task, with the slight self-doubt blip in
Chat being the only deviation.

An aggregated Friedman test across the six subscales is not
significant ($\chi^2(2) = 5.24$, $p = 0.073$, $W = 0.038$), so the
composite workload score does not differ significantly between
modes. Two subscales reach significance individually: Temporal
Demand ($\chi^2(2) = 12.77$, $p = 0.002$, $W = 0.093$) and
Frustration ($\chi^2(2) = 8.93$, $p = 0.012$, $W = 0.065$).
Pairwise Wilcoxon signed-rank tests with Bonferroni correction
confirm Form differs significantly from both Ask and Chat on
these subscales, with medium-to-large rank-biserial effects
($|r| = 0.37$ to $0.57$); Ask versus Chat is not significant on
any measure. Full statistics are in Appendix~\ref{app:inferential}.

Overall, the R-TLX reinforces the SUS Form lead: Form sits below
the LLM-driven modes on every workload dimension except
Performance, with two subscales (Temporal Demand and Frustration)
reaching inferential significance individually even when the
composite test does not. The user-study side's strongest single
inferential signal therefore lands on perceived workload, not
usability. Whether perception of time pressure matches actual
time taken is the question for the supplementary completion-time
analysis that follows.


\paragraph{Supplementary: per-mode completion time.}

The R-TLX workload data is entirely perceptual. As a behavioral
counterpart to Temporal Demand, every ROPA document generated by
ROPAgen carries a \texttt{Time elapsed: $N$ seconds} line at
submission, giving one objective completion-time measurement per
participant per mode on the $n = 113$ document sample. The
question is whether the perceived time pressure (Form lowest,
Chat highest) is mirrored by the actual time spent.
Form takes about half the wall-clock time of either LLM-driven
mode. Median completion time was $7$ minutes for Form ($420$~s),
$14$ minutes for Ask ($858$~s), and $12$ minutes for Chat
($717$~s); the means tell the same story (Form $564.2$~s, Ask
$989.8$~s, Chat $979.1$~s), with substantial right-skewed spread
in the LLM-driven modes (Chat SD $\approx 1015$~s, one extreme
session at $6199$~s). The full distributions are in
Figure~\ref{fig:per-mode-timing}.

\begin{figure}[htbp]
\centering
\includegraphics[width=0.85\linewidth]{per_mode_timing_boxplot.png}
\caption{Per-mode completion time in seconds, extracted from the
\texttt{Time elapsed} line emitted by ROPAgen at document
submission ($n = 113$ documents: Form $37$, Ask $37$, Chat $39$).}
\label{fig:per-mode-timing}
\end{figure}

A Kruskal--Wallis test on the full $n = 113$ sample confirms that
completion time varies by mode ($H = 16.46$, $p < 0.001$,
$\varepsilon^2 = 0.15$, large). Form is significantly faster than
both Ask (median $420$~s vs $858$~s; $r = 0.512$, large) and Chat
(median $420$~s vs $717$~s; $r = 0.418$, medium); Ask and Chat
are statistically indistinguishable from each other.

The behavioral data sharpens the perceived workload picture: Form
is not merely felt as the least time-pressured mode, it is also
measurably the fastest by a margin of roughly seven minutes per
document over either LLM-driven mode. The two LLM-driven modes,
despite their very different interaction styles (per-section
explainer versus freeform conversation), take essentially the
same amount of time on average.
Whether participants nevertheless ascribe more value to the
LLM-driven modes when asked specifically about AI support is the
question taken up next in Section~\ref{sec:eval-ai-support}.
